\documentclass[11pt]{article}
\usepackage[margin=2.3cm]{geometry}
\usepackage{xcolor,enumitem,amsmath}
\usepackage[hidelinks]{hyperref}
\newenvironment{comment}{\par\medskip\noindent\begin{minipage}{\linewidth}\itshape\color{black!70}}{\end{minipage}\par\smallskip}
\newcommand{\dt}{$\Delta T$}
\newcommand{\resp}{\noindent\textbf{Response.} }
\newcommand{\change}{\par\noindent\textbf{Change.} }
\setlength{\parindent}{0pt}\setlength{\parskip}{4pt}

\begin{document}
\begin{center}
{\large\bfseries Response to reviewers}\\[2pt]
SMHL-D-26-00918, ``Accuracy does not buy timing: a lead-time metric for
stress anticipation in mobile and wearable sensing''\\[2pt]
(revised title: ``\ldots{} in mobile sensing''; see Reviewer 2, comment 4)
\end{center}

Dear Editor and Reviewers,

Thank you for the careful reviews. They led us to re-examine the
analysis pipeline end to end, not only the points raised. We summarise
the main changes first, then answer each comment in turn. Section, table and figure numbers refer to the revised manuscript;
in the tracked-changes version, additions are in blue and deletions in
red.

\paragraph{Summary of main changes}
\begin{enumerate}[leftmargin=*,itemsep=1pt]
\item \textbf{Stress trajectory.} \dt{} is now computed on the expected
stress level $s_t=\sum_c c\,p_t^{(c)}$, which represents stress
directly, instead of the maximum class probability, which measures model
certainty (R1-1, R2-2). The maximum-probability results move to
Appendix~A.
\item \textbf{One source for every number.} All StudentLife models were
retrained, and every table and figure is generated by a single script
from one set of records. This resolves the figure/table mismatches
(R1-Q1, R1-Q2).
\item \textbf{New analyses:} null models (R2-1); onset-paired
comparisons with a student-clustered bootstrap (R1-5); a TOST
equivalence test with a stated margin (R1-3); a prediction-horizon
experiment (R2-3); a full penalty $\times$ matching-window sensitivity
grid (R1); exact $p$ values and effect sizes with confidence intervals
(R1).
\item \textbf{Narrower claims.} The equivalence claim is now limited to
the quantity tested (the paired mean \dt{} difference on onsets matched
in both models). The chance-level statements follow the null-model
results exactly. The forecast-trained models are compared with the
present-state models on identical validation sequences.
The paper also states that change points are located retrospectively,
so \dt{} measures temporal alignment and not real-time warning lead.
\item \textbf{Corrections we found ourselves.} (i) The reference
implementation's cross-validation is stratified by student and label at
the example level, not split by subject; the text previously said the
opposite. (ii) \texttt{ruptures} evaluates change points on a grid of
every fifth sample by default; we now set \texttt{jump}~$=1$. (iii)
The reference WESAD pipeline evaluates with a shuffled data loader and
stores no window index, so the WESAD trajectories we analysed were not in
temporal order. We have withdrawn the WESAD \dt{} analysis and removed
``wearable'' from the title.
\end{enumerate}

%──────────────────────────────────────────
\section*{Reviewer 1}

\begin{comment}1. The confidence trajectory does not consistently represent stress.\end{comment}
\resp We agree. The maximum class probability is high whenever the model
is confident, including confidently \emph{unstressed}, so a change point
in it need not be a change in predicted stress. We now use the expected
stress level $s_t$ (Eq.~2), which is monotone in the probability mass on
higher stress classes. Appendix~A reports the submitted max-probability trajectory and
$P(\text{high})$ alongside $s_t$. Under the expected-stress trajectory
the main result holds in the form we now state it: the paired mean
\dt{} differences between models are within the $\pm 6$~h margin
(Section 6.3), with anticipation rates of 22.8--26.6\% and medians of
0~h. Appendix~A is a descriptive sensitivity table; we ran no
equivalence test on the alternative trajectories. It shows the same
pattern under each (medians of 0~h, anticipation rates of 22--27\%).
\change Section 5.1, ``Stress trajectory''; Appendix A; terminology
unified to ``stress trajectory'' throughout.

\begin{comment}2. The asymmetric loss is not mathematically defined
(temporal window, weights, normalisation, onset handling, role of $\lambda$).\end{comment}
\resp We now give the full definition. Onsets are upward label
transitions in each training subject's chronological sequence, computed
from training labels only. Each training example receives weight
$1+\lambda$ if it lies within $H=72$~h before its nearest onset,
$(1+\lambda)^{-1}$ if it lies within $H$ after it, and 1 otherwise
(Eq.~3). The weighted, class-weighted cross-entropy is averaged over the
mini-batch without renormalisation (Eq.~4). The reconstruction term,
validation loss and checkpoint selection are unchanged. $\lambda=0$
recovers the reference objective; the pre/post weight ratio is
$(1+\lambda)^2$. The text also gives the numerical settings (class
weights 0.6456, 0.5635, 1.0000; $\beta = 1$; $\alpha = 10^{-4}$ with an
L1 reconstruction term summed over elements, which does not apply to
the baseline configuration used for this experiment, in which the
reconstruction term is disabled and the decoder is not trained;
mini-batch size 1) and
the tie rule (the earlier of two equally near onsets).
\change Section 6.6, Eqs.~3--4.

\begin{comment}3. The invariance claim needs an equivalence margin and an
equivalence/non-inferiority test (e.g.\ TOST).\end{comment}
\resp Agreed: a non-significant difference is not evidence of
equivalence. We pre-specify a margin of $\pm 6$~h, a quarter of the
daily EMA cycle and the smallest advance we consider actionable for a
same-day intervention. We test equivalence with a TOST on the
onset-paired mean difference, using a student-clustered bootstrap.
Because a reader may prefer another margin, we also report the smallest
margin at which equivalence holds for each pair. The margin was fixed
when the revision analysis was designed, before the models were
retrained. The TOST is computed on the student-level bootstrap
distribution of the paired mean, so it accounts for several onsets
coming from one student. For all three pairs the paired mean difference
is within the $\pm 6$~h margin (TOST $p \leq 0.005$), and in fact
within 4.5~h (baseline vs.\ 1 branch), 3.3~h (baseline vs.\ 3 branches)
and 2.8~h (1 branch vs.\ 3 branches). We have been careful about what
this shows. It bounds the \emph{mean} shift among onsets that both
models respond to. It does not establish that the \dt{} distributions
have the same shape, tails or anticipation rate, and the abstract,
contributions, Section 6.3 and conclusion now say ``paired mean
differences within the margin'' instead of ``equivalent distributions''.
\change New Section 6.3, Table 4, Figure 5.

\begin{comment}4. The cross-modal comparison is uncontrolled and relies on
the anticipation rate, despite the paper's own caution in Section 6.3.\end{comment}
\resp We agree, and have removed the cross-dataset claim. The abstract,
contributions and conclusion no longer state a ``24.8--33.3\% band''.
We went further than relabelling it as descriptive: for the reason
given under Reviewer 2, comment 4, the WESAD \dt{} analysis has been
withdrawn altogether, so the revised paper makes no cross-modal claim
and its title no longer refers to wearable sensing.
\change Title; Abstract; Contributions; Results (WESAD subsection and
the cross-modal table removed); Limitations (``Withdrawn cross-modal
analysis''); Conclusion.

\begin{comment}5. Observations are dependent; use paired comparisons of
matched onsets across models.\end{comment}
\resp The folds depend only on the labels, so all configurations share
identical validation sets and identical onsets. For each onset we take
the nearest change point of each model (median across repetitions) and
compare models on the paired differences. Inference uses a bootstrap
that resamples students, which handles dependence both within subjects
and across repetitions. The paired mean differences are $+2.16$~h
[95\% CI $-0.88$, $+5.10$] for baseline vs.\ 1 branch, $+0.43$~h
[$-3.03$, $+3.73$] for baseline vs.\ 3 branches, and $-0.97$~h
[$-3.14$, $+1.29$] for 1 vs.\ 3 branches (163--173 paired onsets from
22 students). The intervals and the TOST come from the student-level
bootstrap. An onset-level Wilcoxon $p$ value is also shown (0.51, 0.83,
0.09) but is labelled descriptive, because it treats onsets as
independent. Onsets matched in one model but not the other are excluded
from that pair, and the text says so. The pooled Wasserstein
and Mann-Whitney results are kept as secondary descriptions and labelled
as such.
\change Section 6.3, Table 4.

\begin{comment}6. Figures 2, 4 and 6 need formatting.\end{comment}
\resp All figures were regenerated from the new records, with consistent
fonts, no overlapping titles, and vector output. The bar chart of
Wasserstein distances (old Figure 4) is replaced by Figure 5, which
shows the paired differences with confidence intervals against the
equivalence margin. The asymmetric-loss figure (now Figure 7) was
redrawn.

\begin{comment}7. Proofreading: run-in headings in tables, repeated words,
inconsistent terminology, spacing, wrong citations in the introduction.\end{comment}
\resp Corrected throughout. Specifically: run-in paragraph headings were
reformatted; terminology is unified (``stress trajectory'',
``anticipation rate'', ``matched pair'', ``onset''); the introduction now
cites a paroxysmal-AF prediction study for cardiac lead time
(Ebrahimzadeh et al., 2018) and a glucose-forecasting benchmark for
prediction horizons (De Bois et al., 2022). Bibliographic errors were
also fixed (Gjoreski et al.\ is 2017; Jackson et al.\ author list).

\begin{comment}Minor: ``Section ??'' in the WESAD discussion.\end{comment}
\resp The unresolved reference was removed with the withdrawn WESAD
subsection. Every remaining cross-reference was checked in the compiled
PDF.

\begin{comment}Q1: Figure 4 shows $W$ = 4.24/6.50/2.92 h but Section 6.1
reports 3.55/5.75/2.64 h. Q2: Figure 2 shows 24.5/25.5/30.2\% but Table 2
reports 24.8/25.6/30.7\%.\end{comment}
\resp Thank you for catching these. In the submitted version the
figures and the tables were produced by separate runs of the analysis,
and the figure annotations pooled matched pairs across repetitions
while the tables reported repetition means. We cannot reconstruct the
submitted figures exactly, because the model outputs behind them are no
longer available to us, so we did not try to patch the numbers. Instead
we retrained every model and now generate every table and every figure
from one set of records in a single pass, with figure annotations taken
from the same values as the tables. The anticipation rate is reported
both pooled and as a repetition mean $\pm$ SD, each labelled. The
revised values also differ from the submitted ones for a reason we
found during the revision and describe in Section 5.2: the submitted
analysis used the \texttt{ruptures} default candidate grid of every
fifth sample. Appendix~A reports the submitted configuration re-run on
the new models; it reproduces the scale of the submitted numbers
(332--420 matched pairs, medians of $+4.0$ to $+5.5$~h, anticipation
24--25\%).

\begin{comment}Suggestions: exact $p$ values and effect-size CIs instead
of ``$p>0.20$''; full asymmetric-loss maths; full StudentLife penalty and
matching-window sensitivity.\end{comment}
\resp All three were done. Table 6 reports exact Mann-Whitney $p$
(0.783 and 0.940) and, as the effect size, the onset-paired mean
difference from the symmetric baseline with a 95\% student-clustered
bootstrap CI ($-0.13$~h [$-2.73$, $+2.86$] and $-1.78$~h [$-5.26$,
$+0.95$]). The Hodges-Lehmann shift is 0.0~h at both strengths. The
macro-F1 gain is now 3.4 and 3.1 points (fold-paired Wilcoxon
$p = 0.002$ and $0.010$), with a small loss of micro-F1 that we report. The sensitivity grid (penalty
$\in\{0.5,1,2,5\}$ $\times$ window $\in\{24,48,72,96\}$~h) is in
Figure 9, and every cell, including penalty 5, is tabulated in
Appendix~B. Across penalty $\in \{0.5, 1, 2\}$ and window
$\in \{24, 48, 72, 96\}$~h, the anticipation rate stays within
21--34\% (21--28\% at penalty $\leq 1$), rises mechanically with the
window, and the median \dt{} is
0~h in every cell; no cell shows a personalisation advantage in timing.
Penalty 5 leaves fewer than a dozen matched pairs, so those rates are
not meaningful; they are listed in Appendix~B for completeness.

%──────────────────────────────────────────
\section*{Reviewer 2}

\begin{comment}1. (Major) There is no null model: random change points or
shuffled labels.\end{comment}
\resp We added two nulls for each configuration, with 1{,}000 draws
each. (a) \emph{Random change points}: the same number of change points
per subject, placed uniformly at random. (b) \emph{Shifted labels}: each
subject's label sequence is circularly shifted, which preserves its
autocorrelation and onset count while breaking alignment with the model.
We compare the anticipation rate, the median \dt{} and an alignment
statistic (the share of change points within $\pm24$~h of an onset).
No configuration shows significantly more anticipation than either
null. Against random change points (mean $\approx 26$\%), the
personalised models' rates of 26.4\% and 26.6\% are within the null
interval, and the baseline's 22.8\% is significantly \emph{below} it
($p < 0.01$). Against shifted labels (38--40\%), every observed rate is
far lower ($p < 0.01$). On alignment the picture is mixed, and we report
it as such: the baseline's change points fall within $\pm 24$~h of an
onset more often than random change points do (51.0\% vs.\ 46.5\%,
$p < 0.01$) but not more often than under shifted labels ($p = 0.07$),
and neither personalised model is above either null. The paper no
longer says that a quarter of transitions carry a detectable precursor;
a rate near 26\% is what the matching rule gives without timing
information.
\change New Section 6.2, Table 3, Figure 4.

\begin{comment}2. (Major) Max-probability confidence tracks certainty, not
stress; WESAD uses $P(\text{class 1})$, which is inconsistent with
StudentLife.\end{comment}
\resp Agreed; see our response to R1-1. StudentLife now uses a stress
trajectory, and the inconsistency disappears with the withdrawal of the
WESAD analysis (comment 4).

\begin{comment}3. Models predict the current label, so the result is not
surprising. Shift labels forward, retrain, and show how accuracy decays
with horizon.\end{comment}
\resp We accept the premise and ran the experiment. Labels were shifted
1, 2 and 3 EMA responses ahead (median gaps of 16, 30 and 49~h), and
the baseline and three-branch models were retrained (5 repetitions
$\times$ 5 folds each). \dt{} is still measured against current-label
onsets. Three things came out of it.

\emph{Accuracy does not decay with horizon.} Micro-F1 goes from 0.478
to 0.490 for the baseline and from 0.611 to 0.626 for the three-branch
model at the 49-hour gap. Self-reported stress is autocorrelated enough
that forecasting is as easy as nowcasting.

\emph{A direct comparison with the $h = 0$ runs is misleading.} Taken at
face value, the mean \dt{} is 3 to 6~h earlier and the anticipation
rate 4 to 11 points higher in the forecast runs, and the difference in
means is nominally significant in five of six runs. But the forecast
runs are validated on different sequences (the last examples of each
student are dropped and the folds change), and a random-change-point
null on those sequences moves in the same direction with no model
involved.

\emph{On identical sequences the apparent shift is not reproduced.} Every example
has an out-of-fold prediction from the present-state model, so we
rebuilt the present-state model's trajectory on each forecast run's
exact validation sequences and compared the two models onset by onset
(Table 9). The paired mean differences are between $-1.0$ and $+2.0$~h
with no consistent sign, every 95\% interval includes zero, and the
difference is within the $\pm 6$~h margin in five of six runs (the
baseline at $h = 2$ is the exception, with an interval reaching
$+7.1$~h). In summary: on identical validation sequences, no
statistically detectable paired mean timing difference was found.
Equivalence within $\pm 6$~h was supported for five of six comparisons;
the baseline at $h = 2$ remained inconclusive. We also report the lesson this taught us: \dt{}
summaries from different validation sequences are not comparable, and
the Limitations now say so.
We have also rephrased the framing: the paper now states explicitly that
present-state models are not expected to anticipate, and that the
contribution is a metric that makes this measurable and shows that the
standard accuracy metrics do not reveal it.
\change New Section 6.7, Tables 7--9, Figure 8; Limitations (``Dependence on the evaluation sequence''); Introduction, paragraph 4.

\begin{comment}4. WESAD/TSST: the 66 pairs share onsets and are not
independent; median $\Delta T=0$ is uninformative; the rate is 25.0\% at
$\rho=1.0$, and ``stable at 33\%'' cherry-picks three penalties; the
abstract still uses 24.8--33.3\%.\end{comment}
\resp These criticisms are correct, and examining them uncovered a more
basic problem. The reference WESAD pipeline evaluates with a shuffled
data loader and stores no window index, so the trajectories we analysed
were not in temporal order. The submitted WESAD \dt{} values were
therefore not valid timing measurements, and no reweighting of the 66
pairs could repair that. We have withdrawn the WESAD \dt{} analysis.
The paper now rests on StudentLife alone, the title no longer mentions
wearable sensing, the ``24.8--33.3\%'' band is gone from the abstract,
contributions and conclusion, and the Limitations section states what
was withdrawn and why. A cross-modal test with a temporally ordered
evaluation is listed as the first item of future work.
\change Title; Abstract; Sections 1, 3, 5 and 6; Limitations; Future
directions; Conclusion.

\begin{comment}5. Writing: ``Section ??''; overlapping title text in
Figures 2 and 6; Islam et al.\ (2020, HARDC) is miscited as early
arrhythmia warning.\end{comment}
\resp All fixed; see R1-7 and R1-minor. HARDC is a heartbeat
classification method, not an early-warning system, and has been
replaced by Ebrahimzadeh et al.\ (2018) on predicting paroxysmal atrial
fibrillation.

\section*{Further changes not requested by the reviewers}
\begin{itemize}[leftmargin=*,itemsep=1pt]
\item \textbf{Anticipation-rate caution reworded.} The submitted text
said that any widening of the \dt{} distribution raises the
anticipation rate. That is not true in general (it is unchanged for a
distribution symmetric about zero), and Section 6.5 now says ``can''.
The section also no longer calls the rate difference between models a
false positive: the difference in the rate is real; what it does not
show is a shift in the median or the paired mean.
\item \textbf{Retrospective change points.} Section 5 and the
Limitations state that PELT is an offline procedure, so \dt{} measures
retrospective alignment and not the lead time of a real-time alert.
\item \textbf{Macro-F1 mechanism.} The gain under the asymmetric loss is
now reported with the accompanying micro-F1 loss, and class rebalancing
is offered as an interpretation, since per-class scores were not
recorded.
\end{itemize}

\bigskip
Sincerely,\\
David Olutunde Daniel
\end{document}
